%%%%%%%%%%%%%%%%%%%%%%%%%%%%%%%%%%%%%%%%%%%%%%%%%%%%%%%%%%%%%%%%%%%%%%%%%%%%
% AGUJournalTemplate.tex: this template file is for articles formatted with LaTeX
%
% This file includes commands and instructions
% given in the order necessary to produce a final output that will
% satisfy AGU requirements, including customized APA reference formatting.
%
% You may copy this file and give it your
% article name, and enter your text.
%
%
% Step 1: Set the \documentclass
%
%

%% To submit your paper:
\documentclass[draft]{agujournal2019}
\usepackage{url} %this package should fix any errors with URLs in refs.
\usepackage{lineno}
\usepackage[inline]{trackchanges} %for better track changes. finalnew option will compile document with changes incorporated.
\usepackage{soul}
\usepackage{natbib}
\usepackage{amsmath}
\def\permille{\ensuremath{{}^\text{o}\mkern-5mu/\mkern-3mu_\text{oo}}}
\usepackage{lmodern}
\usepackage{float}

% DELETE this part
\usepackage{xargs}
\usepackage[colorinlistoftodos,prependcaption,textsize=tiny]{todonotes}
% 
\newcommandx{\josue}[2][1=]{\todo[linecolor=red,backgroundcolor=red!25,bordercolor=red,#1]{#2}}
\newcommandx{\camille}[2][1=]{\todo[linecolor=blue,backgroundcolor=blue!25,bordercolor=blue,#1]{#2}}
% DELETE up to here

\linenumbers
%%%%%%%
% As of 2018 we recommend use of the TrackChanges package to mark revisions.
% The trackchanges package adds five new LaTeX commands:
%
%  \note[editor]{The note}
%  \annote[editor]{Text to annotate}{The note}
%  \add[editor]{Text to add}
%  \remove[editor]{Text to remove}
%  \change[editor]{Text to remove}{Text to add}
%
% complete documentation is here: http://trackchanges.sourceforge.net/
%%%%%%%

% \draftfalse

%% Enter journal name below.
%% Choose from this list of Journals:
%
% JGR: Atmospheres
% JGR: Biogeosciences
% JGR: Earth Surface
% JGR: Oceans
% JGR: Planets
% JGR: Solid Earth
% JGR: Space Physics
% Global Biogeochemical Cycles
% Geophysical Research Letters
% Paleoceanography and Paleoclimatology
% Radio Science
% Reviews of Geophysics
% Tectonics
% Space Weather
% Water Resources Research
% Geochemistry, Geophysics, Geosystems
% Journal of Advances in Modeling Earth Systems (JAMES)
% Earth's Future
% Earth and Space Science
% Geohealth
%
% ie, \journalname{Water Resources Research}

\journalname{Geophysical Research Letters}

\hfuzz=100pt
\hbadness=10000

\begin{document}

%% ------------------------------------------------------------------------ %%
%  Title
%
% (A title should be specific, informative, and brief. Use
% abbreviations only if they are defined in the abstract. Titles that
% start with general keywords then specific terms are optimized in
% searches)
%
%% ------------------------------------------------------------------------ %%

% Example: \title{This is a test title}

\title{Buoyancy fluxes of sea-ice leads in the Arctic}


%% ------------------------------------------------------------------------ %%
%
%  AUTHORS AND AFFILIATIONS
%
%% ------------------------------------------------------------------------ %%

\authors{Josu\'e Mart\'inez-Moreno\affil{1}, Camille Lique\affil{1}, Claude Talandier\affil{1}, someone else?}


\affiliation{1}{Laboratoire d'Oc\'eanographie Physique et Spatiale (LOPS), University of Brest/IFREMER/IRD/CNRS, Brest, France}

\correspondingauthor{Josu\'e Mart\'inez-Moreno}{josue.martinez.moreno@ifremer.fr}

\begin{keypoints}
    \item Buoyancy fluxes in Linear Kinematic Features vary seasonally 
    \item Buoyancy fluxes in Linear Kinematic Features can be 20-30\% larger than in the ice-pack in summer.
    \item The presence of Linear Kinematic Features has a small impact in the transform water masses at the ocean surface.
\end{keypoints}

\begin{abstract}
    The Arctic sea-ice, in particular the ice-pack, acts as an insulator between the atmosphere and the ocean below. Linear kinematic features (LKF), common in Arctic sea ice, serve as openings that allow exchanges between ocean and atmosphere, for example: heat, salt, moisture, and gases. Therefore, LKFs can influence heat exchanges, ice melt or growth, ocean dynamics, and the atmospheric boundary layer. Existing literature has suggested that LKFs play a significant role in the heat budget of the surface of the ocean and likely impact water mass transformation within the Arctic. However, due to the lack of observations and high resolution models capable to resolve LKFs, an estimate of the impact of LKFs in the ocean has been elusive.  Here, we explore the magnitude and impact in the ocean surface caused by the LKFs using SEDNA, a really high-resolution pan-Arctic hindcast (1/60$^\circ$; NEMO + SI3 configuration). We show that despite the LKFs accounting for a small area of the sea-ice cover (<20\%), they contribute up to 10\% more buoyancy than the ice-pack at the surface. Furthermore, of the total surface water mass transformation in the Arctic ice cover, LKFs explain approximately 10\% of it and it follows the same temporal evolution as in the ice pack. The present estimate shows a significant local impact but a smaller modulation of the Arctic Ocean dynamics as previously hypothesized.
\end{abstract}

\section{Introduction}
% Paragraph topic: Discuss Sea ice, insulator properties, and LKF

Arctic sea-ice plays a crucial role in regulating Earth's climate by acting as a natural insulator between the atmosphere and the ocean below \citep{Wettlaufer_insulator_1997, Untersteiner_ice_insulator_1961}. This insulating property acts as a thermal barrier, limiting heat loss from the ocean to the atmosphere and minimizing evaporation, thereby maintaining the stability of the polar environment \citep{Ruffieux_pack_lead_1995}. A ubiquitous feature of Arctic sea ice is the formation and persistence of linear kinematic features (LKF), commonly known as leads. These features are elongated openings in the sea ice cover that expose the ocean surface to the atmosphere \citep{Lupkes_leads_2008}. 

% Paragraph topic: LKFs scales, LKFs formation mechanisms + climatology.

Leads vary in scale from a few meters to kilometers in width and span distances of a few kilometers up to hundreds of kilometers \citep{Tschudi_obs_leads_1998}. Their formation mechanisms are diverse, ranging from mechanical to thermodynamic processes \citep{Rampal_nextsim_2016, Hutchings_model_lkf_2005}. Mechanically, sea ice deformation and fracturing are caused by winds and ocean currents \citep{Richter_pack_stress_2002}. Winds are thought to be the main forcing of the sea ice deformation, where persistent winds stress result in sea ice fracturing and eventually form LKFs \citep{Linow_lkf_2017, Hutter_scaling_lkf_2018}. Although ocean currents have a smaller contribution to the sea ice deformation compared to wind forcing, ocean drag at the bottom of the ice acts to balance the wind-induced sea ice velocities \citep{Weiss_scale_2004}. Additionally, observations show that ocean currents may play a bigger role in dictating the fracturing patterns of sea ice in regions with strong currents and bathymetric features \citep{Willmes_leads_currents_2023}. Alongside these mechanical forces, thermodynamic processes such as rapid variation of the air or ocean temperature can result in the fracturing of sea ice \citep{Lewis_thermomechanics_1998}. Once a lead is formed, warm air temperatures can instigate localized melting, maintaining the lead open. Conversely, during colder periods, rapid freezing facilitate the closure of a lead. The intricate interplay between these thermodynamic processes and dynamic forces dictates the evolution and persistence of leads.

% \textbf{Describe believed importance in the fluxes and likely impact to WMT.}
Leads are crucial since they act as gateways between the atmosphere and the ocean, exchanging heat, moisture, and gases \citep{Michaelis_lead_2022, Andreas_vapor_2002}. Depending on the season and the atmospheric forcing, this direct contact between the ocean and atmosphere can result in locally intensified heat fluxes at the ocean-atmosphere interface \citep{Marcq_flux_lead_2012,Maykut_heat_1986}, localized ice melt or growth, and modify the ocean dynamics \citep{Bourgault_flux_lkf_2020}. For example, in winter, heat exchange from the ocean to the cold atmosphere, in addition to ice freezing within the leads, results in local convective processes in both the ocean and atmosphere \citep{Lupkes_leads_2008,McPhee_LKF_2005}. A change in the surface temperature and salinity due to the heat and salt flux at the ocean-ice-atmosphere interfaces transforms locally the surface water masses \citep{Pemberton_WMT_2015, Walin_WMT_1982}. Leads have a undeniable local impact on the ocean and it has been hypothesized that leads may play a significant role in the Arctic by modulating the stratification, altering ocean circulation, and transformation of water masses \citep{Reiser_LKF_2020,Smith_polynyas_1990}. However, the magnitude and influence of the LKFs compared to the surrounding ice pack in the Arctic basin is yet to be quantified.

To investigate and quantify the impact of the LKFs in the Arctic Ocean, we use a really high-resolution pan-Arctic hindcast, SEDNA \citep{SEDNA_2023}. Then, we compute a mask of the linear kinematic features during the year 2014, as well as the buoyancy fluxes and water mass transformation for the identified LKFs, and compare them to the ice pack (See Methods for details of the calculations and numerical model).  

\subsection*{Observation vs Models}

\begin{figure}
    \centering
    \includegraphics[width=0.9\textwidth]{./figures/Figure_1_SENTINE_SEDNA_LKF_conc.pdf}
    \caption{Linear kinematic features from SEDNA. a) Linear kinematic mean frequency during winter between 2011 and 2015. b) Snapshot of the total deformation for the 21st of April 2014. c) True color image from Sentinel 3 on the 12th of April 2022 for the magenta area in panel a. d) Mean ice concentration from SEDNA (800m resolution) on the 21st of April 2014. e) Total deformation used to identifies the linear kinematic features. f) Identified linear kinematic features on the 21st of April 2014.}
    \label{fig:fig1}
\end{figure}

% \textbf{Figure 1}
% \begin{itemize}
%     \item Describe spatial pattern of features, for example in satellite we observe many floes of different scales that the model can not reproduce
%     \item Describe that the model is a mean state over a day, thus it is likely smoother. 
%     \item Highlight that model still reproduces features similar to those observed in the model. 
%     \item The length-scales also appear to match (Should we do something to quantify this?)
%     \item Add statistics of the LKFs compare to other studies? 
% \end{itemize}

Climate models and high resolution sea ice models are unable to reproduce LKFs due to the lack of resolution to capture the small footprint of LKFs in the Arctic \citep{Wang_leads_2016}. SEDNA is capable of overcoming this limitation, since its spatial resolution of $\sim 800m$ on average across the Arctic suffice to be reproduce the observed sea ice deformation statistics \citep{Bouchat_sea_ice_rheology_2022}. The climatology of the winter lead probability, i.e. the likelihood of a lead occurring for any given day during winter, is shown in figure \ref{fig:fig1}a. There are several distinct spatial patterns that match those from observations \citep{Willmes_leads_currents_2023}. For example, we observe increased lead frequency in the shelves, along bathymetric features, and in regions with large kinetic energy. Several hotspots can be identified along the sea shelves and in the Chukchi Sea, highlighting the model's ability to reproduce LKFs dynamics. \textbf{Should we something more robust to show that the model does a good job?}

SEDNA is capable of reproducing the overall pattern of the observed winter lead distribution in the Arctic shelves, the Greenland Sea, Barents Sea, Kara Sea and Laptev Sea. However, observations show a prominent signature in the Beaufort Sea that is not captured in our simulation using the same winter definition. Using a winter definition starting two month earlier, allows for a better representation of the LKFs in the Beaufort Sea (See supplementary figure 1). A snapshot of the total deformation is shown in figure \ref{fig:fig1}b, where the local maxima in the deformation correspond to the linear kinematic features identified using the \citet{Hutter_identification_2019} algorithm described in the Methods section. A closer examination, focusing north of Greenland in panels c, d, e, and f of figure \ref{fig:fig1}, reveals similarities between the spatial pattern of typically observed leads on the 12th of April 2022 from a true color satellite image (Fig. \ref{fig:fig1}c) and the ice concentration of SEDNA for the 21st of April 2014. Note that the sea ice concentration in the identified leads shown in figure \ref{fig:fig1}e-f remains higher than 80\% concentration; this is a consequence of the visco-plastic nature of the sea ice model and the daily averaged output that smooths the presence of LKFs in the fields. Furthermore, in our simulation, but in fact, in all numerical models using a visco-plastic rheology, it has been observed that they do not reproduce the typical fracturing angles as in observations \citep{Hutter_SIREx_2022}. Despite this, the important similarities between the observations and model lead probability allow us for the first time to estimate quantitatively the impact LKFs have in the ocean.

In situ observations have shown that LKFs can have up to two orders of magnitude larger solar heat input than the ice covered ocean \citep{Maykut_solar_1995} and modify the variance of the heat fluxes in the Arctic basin \citep{Olason_lead_2021}. Presumably impacting the buoyancy flux and the circulation of the Arctic Ocean. 
Figure \ref{fig:fig2}a and e show snapshots of the buoyancy flux, revealing the typical patterns of the Arctic. In winter (Fig. \ref{fig:fig2}a) we observe an overall negative buoyancy flux, a consequence of the freezing of sea ice and thus brine rejection. Meanwhile, in summer, there is an overall positive buoyancy flux due to the melting and freshening of the Arctic (Fig. \ref{fig:fig2}b).  Overall, in these two maps we observe that inside the ice covered region (magenta contour), the fluxes are spatially homogeneous, except for the marginal ice zone and the LKFs. Regarding the LKFs, we observe that these features have larger buoyancy fluxes than the surroundings ice covered ocean, in agreement with previous studies.

\begin{figure}
    \centering
    \includegraphics[width=1\textwidth]{./figures/Figure_2_buoyancy_flux_maps.pdf}
    \caption{Maps of buoyancy flux and transects of downward heat flux, freshwater heat-equivalent flux, ice thickness, and mixed layer depth for the 1st of January and 25th of June 2014. a) Buoyancy flux in Winter the 1st of January 2014. Transect on the 1st of January 2014 of b) downward heat flux, c) freshwater heat-equivalent flux, d) ice thickness and mixed layer depth. e) Buoyancy flux in Summer the 25th of June 2014. Transect on the 25th of June 2014 of b) downward heat flux, c) freshwater heat-equivalent flux, d) ice thickness (magenta) and buoyancy flux (green). Note that the transect location and orientation changes between the winter and summer dates. Magenta contour in panels a and e correspond to the 0\% concentration contour of sea ice. Vertical blue bars shows the mask of LKFs in the transect.}
    \label{fig:fig2}
\end{figure}

% \textbf{Figure 2}
% \begin{itemize}
%     \item Describe the sign of the buoyancy flux (Negave = denser water). 
%     \item Describe that we can see important differences in the BF between the open ocean and the packed ice. 
%     \item  Note some of the LKFs features observed in the buoyancy flux, for example large negative flux in the East Siverian Sea or Beaufort sea and the positive flux in the center of the Arctic.
% \end{itemize}

In figure \ref{fig:fig2}, we plot two transects on the 1st of January 2014 and 30th of June 2014 showing the heat flux (panels b, f), freshwater heat-equivalent flux (panels c and g), and ice thickness and buoyancy flux (panels d and h). Examining the 1st of January 2014, we observe a small heat flux in the LKFs (Fig.\ref{fig:fig2} b); blue shaded areas) that does not impact the heat flux at the ocean-atmosphere interface because the ocean surface is at freezing point. Consequently, the extra heat flux is used for sea ice growth. In contrast, the salt flux (Fig.\ref{fig:fig2}c) shows two peaks with twice the magnitude of the surrounding ice cover collocated within the identified LKFs. A positive freshwater flux corresponds to a salinification of the ocean surface due to sea ice growth and a negative buoyancy flux that increases the density of the  water at the surface (Fig.\ref{fig:fig2}d). As expected, there is a decrease in ice thickness where the LKFs are identified. On the 25th of June (Fig.\ref{fig:fig2}f), we observe an increase in the heat fluxes within some of the identified LKFs to up to $\approx 25 W/m^2$. Meanwhile, the freshwater flux shows a decrease within the LKFs, indicating a freshening of the surface due to sea ice melt. This change is captured by the positive buoyancy flux, decreasing the density at the surface (Fig.\ref{fig:fig2}g and h).  Similarly to the winter snapshot, we note a decrease in ice thickness of more than $1m$ within the LKFs. It is important to note that some peaks in the buoyancy flux align with the edge of the detected leads, likely due to the enhanced refreezing or melting that occur at the edge of the LKFs. Furthermore, not all the identified LKFs exhibit significant ocean-atmosphere fluxes. For these features, we explored snapshots of the simulation a few days later and noted that the LKFs are starting to open, and thus the ocean and ice still behave as they do in the ice pack.% (\textbf{Check this idea?}). 

As in observations, we detect heat fluxes larger in the LKFs than in the surrounding ice pack in both summer and winter, however, the freshwater heat-equivalent flux due to melting and freezing is significantly larger than the heat fluxes associated with the LKFs. Using daily masks of the LKFs, we explore the evolution of the buoyancy fluxes over the full year of 2014 in the next section.


\subsection*{Impact of LKFs in the ocean buoyancy}

We now consider the spatially averaged impact of LKFs on ocean buoyancy. We start by looking at the percentage area covered by four distinct masks of the sea ice covered Arctic Ocean: i) the transitional ice zone (TIZ; 0-80\% sea ice concentration), ii) the ice pack (80-100\% sea ice concentration), iii) the daily identified LKFs in the ice pack, and iv) the daily identified LKFs in the TIZ. The areas covered by the different masks show a seasonal cycle; during summer the percentage of the ice pack cover area drops from approximately $90\%$ to $50\%$ and the transitional ice zone increases from $10\%$ to approximately $50\%$. On one hand, the percentage of the area covered by LKFs in the ice pack peak in winter to almost 20\% coverage over the ice pack area, since winter storms can break up the ice pack. On the other hand, the LKFs in the TIZ peak in summer, to cover up to 25\% of the ice cover Arctic and 50\% of the TIZ surface area, since low concentrated ice can break more easily. On averages the LKFs in the ice pack and the TIZ never exceeds more than 20\% of the sea ice covered Arctic, consistent with previous model assessments \citep{Wang_leads_2016}  

To demonstrate the role of leads in buoyancy flux, we multiply the buoyancy flux and area weighted them by the respective masks of the sea ice cover; the TIZ, the ice pack, the LKFs in the ice pack, and the LKFs in the TIZ. The buoyancy flux shows an important seasonal cycle, where positive values correspond to melting season of sea ice, thus forming lighter waters in the ocean surface, and negative values correspond to the freezing season of sea ice, thus forming denser waters. The buoyancy flux in the TIZ forms generally lighter waters (Fig.\ref{fig:fig3}b), this suggests that in the TIZ there is a preferential melting of sea ice, where both the haline and thermal components of the buoyancy flux act to lighten the ocean surface (Fig.\ref{fig:fig3}c and d). Over the seasonal cycle, the ice pack changes sign from winter to summer (Fig.\ref{fig:fig3}b), this sign change corresponds to a shift from ice formation in winter (formation of denser waters) to ice melting in summer (formation of lighter waters), where the primary contribution to the  buoyancy flux arises from the haline buoyancy component. Similarly the ice pack and TIZ LKFs follow the seasonal cycle of each one of these regions, respectively. The LKFs in the TIZ follows intimately the seasonal cycle, suggesting that the LKFs behave as the surrounding low concentration sea ice. Meanwhile, LKFs in the ice pack follow the seasonal cycle tightly in winter, however in summer the LKFs in the ice pack have larger buoyancy fluxes. The mean summer buoyancy flux of LKFs in the ice pack is 7.26$\times10^{-14} m^2/s^3$, compared to the 5.92$\times10^{-14} m^2/s^3$ in the ice pack. This is approximately 20\% more buoyancy flux in summer suggesting that in summer the LKFs in the ice pack act to enhance melting of the sea ice, likely due to changes in the albedo.

% Analogous to the ice pack, the buoyancy in the LKFs has a seasonal cycle and it shows a similar temporal evolution as the ice pack. Although LKFs cover less than 5\% of the total Arctic ice cover area, over certain periods of time the LKFs buoyancy flux can be 10\%-20\% larger than that of the ice pack.

\begin{figure}[t]
    \centering
    \includegraphics[width=0.9\textwidth]{./figures/Figure_3_buoyancy_flux_components_ts.pdf}
    \caption{Time-series during 2014 of a) the area cover, b) buoyancy fluxes,  c) haline buoyancy flux component, and d) thermal buoyancy flux component  for linear kinematic features in the ice pack, the linear kinematic features in the transitional ice zone, the transitional ice zone (ice concentration 0-80\%), and the ice pack (ice concentration $>$80\%).}
    \label{fig:fig3}
\end{figure}

% \textbf{Figure 3}
% \begin{itemize}
%     \item TODO: Add a time-series of the ratio between the BF-pack/BF-LKF
%     \item Describe that the Buoyancy flux high frequency variability.
%     \item Important seasonal cycle, buoyant water in Summer, denser in Winter.
%     \item Describe that area in LKFs is really small $< 5 %$. 
%     \item Fluxes are mostly dominated by haline component
%     \item Fluxes in LKFs are comparable to those in the Pack ice.
%     \item Explain that despite the area being just a fraction, generally less than 10\% of the ice-covered Artic, the fluxes in LKFs explain the fluxes of the packed ice and their magnitude is comparable to the marginal ice zone. 
% \end{itemize}

\textbf{TODO: Update to include the 4 different masks}
To further quantify the role of the ice pack and the LKFs on the overturning circulation of the Arctic, we use the water mass transformation framework (WMT; \citealp{Walin_WMT_1982}). This approach gives us information on the formation of lighter or denser water due to changes in the surface fluxes. In agreement with previous studies, the surface freshwater flux dominates Arctic water mass transformation. The yearly mean WMT transformation of the ice pack (Fig. \ref{fig:fig4}a) is characterized by a broad region of strong positive transformation (losing buoyancy) in the density range of $22.5 < \gamma_n <27.5$, reaching a maximum of 1.5 Sv, and a negative transformation (gaining buoyancy) below $\gamma_n < 22.5$. The yearly mean WMT in LKFs also  has a positive peak within the $22.5 < \gamma_n <27.5$ reaching up to 0.15 Sv, in other words 10\% of the yearly water mass transformation in the ice covered Arctic. 

The daily WMT of the ice pack and LKFs are shown in figure \ref{fig:fig4}b and c respectively. Note that the magnitude of the WMT in the LKFs is 10\% of that in the ice pack. In both we observe a seasonal cycle. In the winter months, a positive transformation occurs meaning a loss of buoyancy, due to brine rejection due to sea ice growth. In the transition between summer and winter, we reach almost zero transformation due to the fact that both denser and lighter waters are been formed as ice freezes and melts. In summer, a negative transformation or a gain of buoyancy occurs due to the melting of sea ice.  
More importantly, the WMT in the ice pack and the LKFs show a similar seasonal behavior, where some small differences can be observed in the WMT ratio between the LKFs and ice pack. The major differences in the density classes can be observed in $\gamma_n < 22.5$ and $\gamma_n > 27.5$, although this signal is likely to arise in LKFs due to their capacity to form fresher water and denser water when the ocean is in direct contact with the atmosphere, the magnitude of the WMT in these density classes are negligible. 

\begin{figure}[t]
    \centering
    \includegraphics[width=0.9\textwidth]{./figures/Figure_4_water_mass_transformation_hovmoller.pdf}
    \caption{Mean and daily Water mass transformation over 2014. a) Mean water mass transformation over 2014 for the ice pack and LKF. Hovmoller diagram of the water mass transformation (Sv) for the b) Pack Ice and c) linear kinematic features. Pane d) is the ratio of water mass transformation between the LKFs and the Pack Ice. Note the magnitude change in the colorbar of panels a), b), and c).}
    \label{fig:fig4}
\end{figure}

% \textbf{Figure 4}
% \begin{itemize}
% %    \item TODO: Add units below colorbar.
% %    \item TODO: Add mean over the year as a top panel.
%     \item TODO: Map the mean transformation in SV based on the mean density of the Arctic? 
%     \item Describe seasonal cycle of water mass transformation
% %    \item Highlight the change in the magnitude 
% %    \item The LKFs explain around 10\% of the Pack Ice WMT. 
% %    \item During summer the transformation is towards buoyant water and in winter the transformation is towards denser water. 
% %    \item Describe shift of water mass transformation during summer. 
%     \item Is it possible to estimate the effect of water masses transformation based on the outcrop of density contours in the Arctic? 
% \end{itemize}

% \textbf{Should we do a schematic (Fig 1) as in Michaelis, J. \& Lüpkes, C. but for the ocean?}

\section{Conclusions}

Linear kinematic features are ubiquitous in the sea ice covered oceans that allow direct exchanges between the ocean and the atmosphere, whereby fresher and denser water can be formed at the ocean surface depending on the atmospheric forcing. Generally, the importance of these fluxes at the ocean-atmosphere interface have been speculated to be significantly important  for the Artic dynamics \citep{Maykut_solar_1995}. The present study assesses the impact of LKFs in the ocean surface by diagnosing the buoyancy fluxes using high-resolution hindcast (SEDNA). SEDNA reproduces the LKFs spatial distribution over most of the Arctic, with exception of the Beaufort Sea. This mismatch is caused by a two-month delay in the simulated LKFs. In other words, there is a temporal shift in the simulated LKFs frequency and those of observed. This mismatch could have important implications in the Arctic circulation, thus further research is needed to better understand this temporal shift in the formation of leads. %with a brittle Bingham–Maxwell rheology instead of an Visco-Plastic rheology.

% Add paragraph on LKFs in the ice pack being more important in summer. 

Although LKFs have a significant local impact, they contribute only a small fraction to the overall buoyancy flux and water mass transformation of the Arctic. Despite their observable impact on the direct exchanges between the ocean and the atmosphere, LKFs make up only a fraction of the total sea ice cover in the Arctic. Similarly, in terms of water mass transformation, LKFs account for a limited portion of the overall transformation process. While their presence can lead to localized changes in the density of surface waters due to brine rejection and ice melting, their spatial extent and duration are not sufficient to significantly alter the ocean dynamics of the Arctic. Our findings show that LKFs cover a small area of the ice-covered Arctic and they have a contribution of up to 10\% of the buoyancy flux and water mass transformation. 

LKFs generated through shear have been shown to have a significant impact on ocean surface dynamics. For example, \citet{McPhee_LKF_2005} observed a change in the heat fluxes in the ocean subsurface up to two orders of magnitude larger compared to the surrounding areas. Although our definition to identify LKFs include this process, we do not observe such signatures in our simulation. Therefore, further research is required to better understand different LKFs generation mechanisms and their potential impact in the Arctic Ocean. 

Our results may also be sensitive to the specific sea ice and vertical mixing parameterizations used in our model, particularly those relating to ice/ocean fluxes, such as brine rejection, enhanced vertical diffusivity, and vertical exchanges. Finally, it remains to be shown whether LKFs can locally impact regions of the Arctic basin where the sea ice and stratification may allow for a more significant impact on the ocean below LKFs.

% \textbf{reference \citep{Bouchat_sea_ice_rheology_2022} to argue that it's ok to use this high resolution simulation since they seam to resolve properly the LKF.}

% \begin{itemize}
%     \item Describe the SEDNA reproduces well leads.
%     \item Talk about the main results, the small contribution in the large scale impact. 
%     \item Paragraph on LKFs may be important only in the diverging LKF
%     \item Talk about potentially some of the important local impacts of LKF
% \end{itemize}

\textbf{Camille, I'm still working in the Methods, so no need to revise them right now.}

\section{Methods}
\subsection*{Model description}

Here we employ the NEMO numerical platform (release 4.0.5) to establish a 1/60$^\circ$ pan-Arctic ocean-sea ice configuration named SEDNA. The ocean core (NEMO-OCE) solves the primitive equations (PE) using finite differences on an Arakawa C-grid mesh, while the ocean is coupled to the SI3 sea-ice component (NEMO-SI3). The pan-Arctic domain covers the Bering Strait on the Pacific side and extends southward to 56$^\circ$N in the Subpolar North Atlantic gyre and to the Baltic Sea entrance. The ocean model incorporates a linear free surface, utilizes a 3rd order flux form scheme for momentum advection, and employs a 4th order Flux Corrected Transport (FCT) for tracer advection. Additionally, the sea-ice representation includes a 5-category model with the Elasto-Visco-Plastic (EVP) rheology activated and a land-fast capability to better simulate sea-ice behavior in specific areas.

The NCAR bulk formula from Large \& Yeager (2009) to calculate air-sea fluxes and uses hourly ERA5 atmospheric data for near-surface variables. The three lateral boundaries are constrained with daily GLORYS12V1 Re-analysis data, and monthly freshwater fluxes from river discharges are combined with Greenland land ice melt data to supply SEDNA. The simulation extends over seven years (2009 to 2015) and starts from rest, utilizing initial conditions based on the World Ocean Atlas 2009 temperature/salinity and mean January 2009 ice state from PIOMAS re-analysis. Preliminary assessment against available observations and re-analyses data shows that the simulation reasonably captures key ocean and sea-ice characteristics.


\subsection*{Identification of LKFs}

The LKFs from SEDNA were identified by using the algorithm proposed by \citet{Hutter_identification_2019}. This method uses the total deformation rate defined as: 
\begin{equation}
    T_d = \left[\left(\underbrace{\frac{\partial u}{\partial x} +\frac{\partial v}{\partial y}}_{Divergence}\right)^2 + \underbrace{\left(\frac{\partial u}{\partial x} - \frac{\partial v}{\partial y}\right)^2+ \left(\frac{\partial u}{\partial y} + \frac{\partial v}{\partial x}\right)^2}_{Shear}\right]^{1/2},
\end{equation}
where $u$ and $v$ are the sea ice velocities. The deformation rate varies spatially depending on the background deformation and the sea ice properties \citep{}, where the largest deformation rates occur along the boundaries of ice flows where the LKFs are found. Therefore, to identify the LKFs it is required to identify the local maxima using edge detection (DoG filter) in the deformation rate field. After the edge detection, using binary map we obtain a mask with the original width of identified LKF. This mask is then used to compute the statistics presented in this study. After this step, \citet{Hutter_identification_2019} reduce the LKFs width to a single pixel to then track the LKFs over time. Since our study focus on the ocean-atmosphere interactions, we only use the identification algorithm to produce a mask of LKFs during 2014 for the Arctic, where we retain the original width computed from the deformation rate. This mask is then used to quantify and contrast the impact of LKFs compared to the ice-pack.

\subsection*{Buoyancy flux and water mass transformation}
 
% Our LKFs detection algorithm consists of three parts: (1) a preprocessing step that transforms the input deformation data into a binary map of pixels that mark LKFs by using different filter steps, (2) a detection routine that splits the network of LKFs pixels in the binary map into the smallest possible segments, and (3) a reconnection instance that estimates the probability of different segments belonging to one feature and then connects all segments of a LKF. The general structure of the algorithm follows Linow and Dierking (2017), although individual parts have been modified substantially. The main enhancements of the algorithm are a parallel detection of segments with a stronger constraint on the curvature and the introduction of a probability-based reconnection. Further, the entire algorithm was rewritten in Python (Python Software Foundation, http://www.python.org, last access: 15 February 2019) to avoid license issues with the previous code that was based on the commercial software IDL.

\begin{itemize}
    \item Describe MLD definition
    \item Describe buoyancy flux and decomposition between thermal and haline components.
\end{itemize}

\section{Open Research}

\acknowledgments


\bibliography{biblio}

\end{document}